\section{Protected ideals in short flow and transportation games}\label{flow:section}
This appendix concerns a different cooperative technology from the
pair-additive assignment model. It tests which protection statements extend
to arc-player flow games. The results are pointwise statements about each
original entire core. No incentive guarantee across capacity reports is
asserted. In particular, the counterexample below refutes simultaneous
endpoint-ideal attainment while having zero protected subsidy itself.

\subsection{The model and the sharp unit-capacity bound}
Let $G$ be a finite directed acyclic network with distinct source $s$ and
sink $t$. Each arc is owned by a distinct player; parallel arcs are
permitted. A coalition's worth is its maximum $s$--$t$ flow using only its
arcs. Write $T=v(N)>0$, let $\core$ be the entire original core, and set
$H_e=\max_{x\in\core}x_e$. For this appendix define
\[
 s_{1/3}(G)=\min\{y(N)-T:y\ge H/3,\ y(A)\ge v(A)
                                   \text{ for every }A\subseteq N\}.
\]
For any minimum cut, paying each of its arcs its capacity and all other
arcs zero gives a core allocation: its total is $T$, and each coalition's
flow is at most the capacity of its available arcs in that cut. This
also proves nonemptiness of the original core. The core is compact, so
all the maxima defining $H$ are attained.

\begin{lemma}[Deletion and unit ideals]\label{flow:ideals}
For every grand maximum flow $f$ and arc $e$,
\[
                    H_e\le T-v(N\setminus\{e\})\le f_e.
\]
If all capacities are one, $H_e\in\{0,1\}$, and $H_e=1$ precisely
when $e$ belongs to some minimum cut.
\end{lemma}
\begin{proof}
The first inequality follows by combining core coverage of
$N\setminus\{e\}$ with efficiency. Decompose $f$ into paths and remove
all paths using $e$; the lost value is $f_e$, giving the second inequality.
With unit capacities, integral maximum flows show that the deletion
marginal is zero or one. If it is one, a minimum cut after deleting $e$
has capacity $T-1$. The restored arc must cross that cut and makes it a
minimum cut of the original network. Its cut allocation attains payoff
one at $e$. Conversely, membership in a minimum cut gives that same
attainment. If the deletion marginal is zero, so is $H_e$.
\end{proof}
The unit-ideal characterization, cut allocations, and minimum-cut lattice
are classical flow-game ingredients~\cite{flow:Fang2010,flow:KalaiZemel1980}.
They are included to separate them from the protected-subsidy conclusion.

\begin{theorem}[Short unit-capacity flows]\label{flow:unit}
Suppose every directed $s$--$t$ path has at most three arcs, and all arc
capacities are one. Then
\begin{equation}\label{flow:unit-bound}
                     s_{1/3}(G)\le\frac{T-1}{3}.
\end{equation}
For every positive integer $T$, equality occurs in a network whose
$s$--$t$ paths all have three arcs. Consequently the supremum of
$s_{1/3}(G)/T$ in this class is $1/3$.
\end{theorem}
\begin{proof}[Upper bound]
Directed cut capacity is submodular in the source-side vertex set.
The intersection and union of two minimum-cut source sides are therefore
also minimum-cut source sides: both crossed cuts have capacity at least
$T$, while submodularity bounds their sum by $2T$. Since the network is
finite, the intersection $U_-$ and union $U_+$ of all minimum-cut source
sides define minimum cuts $C_-$ and $C_+$.

Every source arc $(s,u)$ with ideal one crosses some minimum cut, so $u$
is missing from at least one source side and hence from $U_-$. Thus
$C_-$ simultaneously contains every positive-ideal source arc. Similarly,
$C_+$ contains every positive-ideal sink arc $(u,t)$. Direct $s$--$t$
arcs cause no difficulty and can be in both cuts.

Let $B$ be the positive-ideal arcs outside $C_-\cup C_+$. Such arcs are
strictly internal. Fix an integral grand maximum flow decomposed into
$T$ arc-disjoint paths. Lemma~\ref{flow:ideals} implies that every arc
in $B$ is used by this flow. Each path has at most one strictly internal
arc, so $r:=|B|\le T$. If $r=0$, the average of the two cut allocations
already dominates $H/3$ and has zero subsidy. If $r>0$, choose $e\in B$
and a minimum cut $C_e$ containing it. Average the allocations of
$C_-,C_+,C_e$, and raise any remaining coordinates to their floors $H/3$.
Only the at most $r-1$ positive-ideal arcs outside these three cuts can
need a raise, of at most $1/3$ each. Raising coordinates preserves every
coalition inequality. This proves~\eqref{flow:unit-bound}.
\end{proof}

\begin{proposition}[A sharp chain family]\label{flow:chain}
For $n\ge1$, use vertices $s,L_1,\ldots,L_n,R_1,\ldots,R_n,t$ and
unit-capacity arcs
\[
 a_i=sL_i,\quad b_i=L_iR_i,\quad c_i=R_it\quad(1\le i\le n),
 \qquad d_i=L_iR_{i+1}\quad(1\le i<n).
\]
Then $T=n$, the ideals of the $3n$ arcs $a_i,b_i,c_i$ are one, all
$d_i$ ideals are zero, and $s_{1/3}(G)=(n-1)/3$.
\end{proposition}
\begin{proof}
The diagonal paths attain flow $n$, matching the source cut. Their
maximum flow avoids the $d_i$, so Lemma~\ref{flow:ideals} gives
$H_{d_i}=0$. The source and sink cuts attain all $a_i$ and $c_i$ ideals.
For $b_i$, the vertex set
\[
           \{s,L_i,\ldots,L_n,R_{i+1},\ldots,R_n\}
\]
has boundary $\{a_j:j<i\}\cup\{b_i\}\cup\{c_j:j>i\}$, a minimum
cut of size $n$. Hence all $b_i$ ideals are also one.

The $n-1$ cross paths $(a_i,d_i,c_{i+1})$ are arc-disjoint. Their union
has worth $n-1$ and leaves $n+2$ positive-ideal arcs outside it. Its
coalition inequality and those exterior floors force protected total at
least $n-1+(n+2)/3$, giving subsidy at least $(n-1)/3$.
For the reverse bound pay $1/3$ to every $a_i,b_i,c_i$, an additional
$1/3$ to $a_i$ for $i<n$, and zero to $d_i$. Every diagonal or cross
path receives at least one. Every coalition's unit-capacity integral
maximum flow decomposes into arc-disjoint such paths, so summing path
inequalities covers every coalition. The total is $n+(n-1)/3$.
\end{proof}
This proves the equality and supremum statements in
Theorem~\ref{flow:unit}. The direct path-cover argument just used depends
on unit capacities; it is not a proof for arbitrary capacities.

\subsection{Prerequisite transfers and a one-sided quota extension}
\begin{lemma}[Prerequisite transfer]\label{flow:transfer}
Let $v$ be a nonnegative monotone game with a nonempty nonnegative core.
Suppose distinct players $a,e$ satisfy
$v(A)=v(A\setminus\{e\})$ whenever $a\notin A$. Transferring all of
$e$'s payoff to $a$ preserves core membership.
\end{lemma}
\begin{proof}
Efficiency and nonnegativity are preserved. Coalitions containing both
or neither player have unchanged payoff, and coalitions containing $a$
only gain. For a coalition containing $e$ but not $a$, the new payoff is
$x(A\setminus\{e\})\ge v(A\setminus\{e\})=v(A)$.
\end{proof}

In a pure three-layer network $s\to L\to R\to t$, let $M$ be the
middle-arc players and $D$ the endpoint-arc players. Define the endpoint
game $w(Q)=v(Q\cup M)$ for $Q\subseteq D$. Its entire core describes
exactly the zero-middle face:
\begin{equation}\label{flow:face}
 \core\cap\{x_M=0\}=\{(z,0_M):z\in\core(w)\}.
\end{equation}
To see this, restrict by testing coalitions $Q\cup M$; conversely, for
any original coalition $A$ use
$v(A)\le w(A\cap D)$. Grand values agree. If every $R_j$ has a unique
sink arc, each middle arc ending there requires that sink arc on any
productive path. Lemma~\ref{flow:transfer} moves all middle payments
sinkward, leaves every source payoff unchanged, and weakly increases
every sink payoff. Thus the source projections of the full core and
this face coincide, and every endpoint coordinate ideal is preserved.
Unique source arcs give the symmetric conclusion by sourceward transfer.
Neither statement changes any original owner or capacity. If both
endpoint ports have parallel owners, neither map is automatically valid.

\begin{theorem}[One-sided quotas]\label{flow:quota}
Fix a rational $q>0$. Consider a pure three-layer network with one or more source
arcs into each $L$ vertex, source capacities $q r_i$ for integers
$r_i\ge0$, exactly one sink arc of capacity $q$ out of each $R$ vertex,
and each rational middle capacity either zero or at least $q$. Parallel source
and middle arcs are allowed; each arc remains a distinct player.
All source ideals are jointly attained in the original core, and all
sink ideals are jointly attained at another core point. If $m$ middle
arcs have positive ideal, then
\begin{equation}\label{flow:quota-bound}
 s_{1/3}(G)\le\frac q3\max\{0,m-1\}\le\frac{T-q}{3}.
\end{equation}
This last bound is attained for every positive integer $T/q$.
\end{theorem}
\begin{proof}
Scale to $q=1$. Because the only arc out of each $R$ vertex has capacity
one, capping every positive middle capacity at one preserves every
coalition value. Zero-capacity arcs have zero core payoff. The normalized
network is integral, so $T$ is a positive integer. The sinkward transfer
and~\eqref{flow:face} preserve all endpoint ideals.

With all middle arcs available, treat each source arc as a firm with
quota $r_i$ and each sink arc as a unit worker. Their assignment value
$a_{ij}$ is one if a positive middle arc links their endpoints, and zero
otherwise. Integral flows give matchings of this many-to-one market.
Conversely, each allowed firm--worker match can be routed through an
available middle arc; a worker is used at most once, so no middle arc
of capacity one restricts the matching. Distinct source arcs sharing a
head are distinct firms with the same allowable workers, not clones
introduced by the proof.

For completeness, we give the classical assignment-game salary-lattice
argument behind side-optimality; a recent reference for this geometry is
\cite{flow:AtayNunezSolymosi2026}. It works for any
nonnegative assignment matrix $a$. Remove zero-quota firms, and balance
total firm quota and worker count by adding zero-value dummy workers
or a zero-value dummy firm. These are null players; their core payoffs
are zero and projecting them away preserves the original core. Complete
an optimal matching $\mu$ to fill every quota and assign every worker.
The added matches have value zero. Let $\mu(i)$ be the workers matched
to firm $i$, and $i(j)$ worker $j$'s matched firm.

A payoff vector $(x,y)$ is in the balanced core precisely when
\begin{align}
 0&\le y_j\le a_{i(j),j},\label{flow:salary-box}\\
 y_k-y_j&\ge a_{ik}-a_{ij}
         &&(j\in\mu(i),\ k\notin\mu(i)),\label{flow:salary-diff}\\
 x_i&=\sum_{j\in\mu(i)}(a_{ij}-y_j).\label{flow:salary-firm}
\end{align}
Indeed, the disjoint coalitions $\{i\}\cup\mu(i)$ attain the grand
value, so their core inequalities are all tight, giving
\eqref{flow:salary-firm}. Deleting one matched worker gives the upper
bound in~\eqref{flow:salary-box}; replacing that worker by an unmatched
worker gives~\eqref{flow:salary-diff}. These use feasible matching
lower bounds, without asserting that a replacement is itself optimal.
Conversely, for any firm $i$ and worker set $U$ of size at most $r_i$,
pair each $k\in U\setminus\mu(i)$ with a distinct
$j\in\mu(i)\setminus U$. Equation~\eqref{flow:salary-diff} gives
$a_{ik}-y_k\le a_{ij}-y_j$, and all remaining matched surpluses are
nonnegative by~\eqref{flow:salary-box}. Hence
\[
                   x_i\ge\sum_{k\in U}(a_{ik}-y_k).
\]
Every feasible coalition matching partitions into disjoint one-firm
teams; summing these inequalities, together with nonnegative unmatched
payoffs, proves all core constraints. Efficiency follows from
\eqref{flow:salary-firm}.

The nonempty compact salary set defined by
\eqref{flow:salary-box}--\eqref{flow:salary-diff} is closed under both
coordinatewise minimum and maximum. For example, two inequalities
$y_k\ge y_j+d$ and $z_k\ge z_j+d$ imply the corresponding minimum
and maximum inequalities. Taking the minimum of finitely many
coordinate-minimizing salary vectors therefore minimizes every salary
simultaneously; by~\eqref{flow:salary-firm} it maximizes every firm's
payoff simultaneously. Coordinate maxima give a point maximizing every
worker payoff. Project away null players and lift through
\eqref{flow:face} to obtain the two asserted original-core endpoint
allocations.

For a normalized middle arc, deletion lowers the integral maximum flow
by zero or one. If it lowers it by one, a minimum cut after deletion,
with the arc restored, is an original minimum cut attaining payoff one
there. Thus each middle ideal is zero or one. Every positive-ideal
middle arc is used by every integral grand maximum flow. Each unit of
flow uses one middle arc, so $m\le T$.

If $m=0$, average the two side-optimal endpoint allocations; this already
dominates $H/3$. If $m>0$, average these two core points and a core point
attaining one positive middle ideal, then raise the other middle
coordinates to their floors. At most $m-1$ coordinates require a raise
of at most $1/3$ each. Rescale by $q$ to obtain~\eqref{flow:quota-bound}.
The chain in Proposition~\ref{flow:chain}, scaled by $q$, lies in this
class and proves sharpness.
\end{proof}
The proof allows parallel source owners because the unique sink supplies
the prerequisite. Arbitrary capacities on both endpoint sides are not
covered: then an endpoint owner may serve several opposite owners, and
the disjoint one-firm-team salary argument need not apply. Splitting
that owner into independent players would change the entire core.

\subsection{An exact transportation embedding}
\begin{proposition}[Embedding without cloning endpoint owners]\label{flow:embedding}
Let disjoint endpoint sets $I,J$ have capacities $r_i,c_j\ge0$ and
allowed pairs $E\subseteq I\times J$. A coalition's worth $w$ is the
maximum transported amount on its allowed pairs, subject to these row
and column capacities, with no separate pair-capacity restriction.
Assume $T=w(I\cup J)>0$. There is a three-layer arc-player flow game
whose entire core is exactly
\[
                  \{(x,0_D):x\in\core(w)\},
\]
where endpoint players retain their identities and $D$ consists of new
middle-arc owners.
\end{proposition}
\begin{proof}
Set $M=\min\{\sum_i r_i,\sum_j c_j\}$. Use source arc $sL_i$ of
capacity $r_i$, sink arc $R_jt$ of capacity $c_j$, and for every allowed
pair two parallel middle arcs $L_iR_j$, each of capacity $M$ and with
its own new owner. With every middle arc available, their capacities
cannot restrict any transportation plan, so endpoint coalitions with
all middle owners have worth $w$.
Deleting any one middle arc leaves its twin of capacity $M$, and hence
does not change the grand value. Every such arc therefore has payoff
zero at every original core point by the deletion bound. Testing endpoint
coalitions together with all middle owners shows that the endpoint
restriction belongs to $\core(w)$. Conversely, extend any point of
$\core(w)$ by zero. The value of any flow coalition is at most the
transportation value of its endpoint restriction, so every original
coalition is covered. Efficiency gives the claimed equality of cores.
\end{proof}
These new middle owners need not be null players of the characteristic
function: they can help coalitions lacking their twins. Only their
zero payoff in every grand-core allocation is used. All returns per
transported unit are one; weighted transportation and a one-unit
constraint per pair are different models.

\subsection{Endpoint ideals need not be jointly attainable}
\begin{theorem}[Twelve-owner endpoint counterexample]\label{flow:endpoint}
In the two-hop network $s\to u\to t$, give the five source arcs
$P,A,G,M,B$ capacities
\[
                         (2,12,46,62,75)
\]
and the seven sink arcs $Q,R,D,E,F,L,N$ capacities
\[
                         (6,10,16,22,24,56,62).
\]
Every arc has a separate owner. Then $T=196$, and the entire-core ideals
are
\[
 H_{\rm source}=(1,6,23,31,38),\qquad
 H_{\rm sink}=(6,10,16,22,24,56,62).
\]
The largest total source payoff is $98$, although the source ideals sum
to $99$. Already $\max_{x\in\core}(x_A+x_B)=43< H_A+H_B=44$.
Nevertheless $s_{1/3}(G)=0$.
\end{theorem}
\begin{proof}
The source capacity sum is $197$ and the sink capacity sum is $196$.
For source set $S$ and sink set $R$, the coalition value is
$\min\{c(S),c(R)\}$. Introduce baseline weights
\[
 a=(1,6,23,31,37),\qquad b=(3,5,8,11,12,28,31),
                  \qquad a(\mathrm{sources})=b(\mathrm{sinks})=98.
\]
Source capacities are $2a$, except that $B$ has one extra unit; sink
capacities are $2b$. We will prove that
\begin{equation}\label{flow:UV}
         U=(a;b),\qquad V=(1,5,23,31,38;b)
\end{equation}
are core allocations, and then prove the matching real-valued upper
bounds on all individual and joint source payoffs.

Let $C_0$ denote the ten players other than $A,B$. Their baseline
weights are $(1,3,5,8,11,12,23,28,31,31)$, of total $153$.
They have no subset of weight $61$, and therefore none of weight $92$.
Here is a direct check. The smaller weights
$D_0=(1,3,5,8,11,12)$, of total $40$, have no subset of weight
$2,7,10$: for $2$ only $1$ fits; for $7$, the weights $1,3,5$ do
not suffice exactly; for $10$, using $8$ leaves $2$, while the smaller
weights total $9$. A subset of weight $61$ uses at most one $31$.
With a $31$, the remaining $30$ either uses $28$ and leaves $2$, uses
$23$ and leaves $7$, or uses only $D_0$ whose complement would have
weight $10$. Without a $31$, using both $28,23$ leaves $10$; using
only $28$ leaves a $D_0$ subset of weight $33$ with complement $7$;
using only $23$ leaves weight $38$ with complement $2$; using neither
cannot reach $61$. These exhaust the cases. Consequently every subset
of all twelve baseline weights of sum $98$ contains $A,B$ either both
or neither: including only $A$ would require a $C_0$ subset of weight
$92$, and including only $B$ would require one of weight $61$.

For arbitrary $S,R$, put $d=a(S)-b(R)\in\mathbb Z$ and
$\delta=\mathbf1_{B\in S}$. If $d\le-1$, coalition worth is
$2a(S)+\delta$, and $U$'s slack is $-d-\delta\ge0$.
If $d\ge0$, worth is $2b(R)$ and its slack is $d\ge0$.
For $V$, these slacks become $-d-\mathbf1_{A\in S}$ when $d\le-1$
and $d+\mathbf1_{B\in S}-\mathbf1_{A\in S}$ when $d\ge0$.
They are nonnegative unless possibly $d=0$. In that case the subset
$S\cup(\mathrm{sinks}\setminus R)$ has baseline weight $98$,
so $A,B$ have equal membership and the slack is zero.
Thus both nonnegative total-$196$ vectors in~\eqref{flow:UV} cover
every coalition.

Now let $x$ be any original core point. For sources define loss
$z_i=c_i-x_i$, and for sinks put $t_j=x_j$. Efficiency yields
\begin{equation}\label{flow:loss-total}
                       \sum_i z_i-\sum_j t_j=1.
\end{equation}
The complement constraint for source $P$ gives $z_P\ge1$;
all source losses are nonnegative by their complement constraints.
Let $\theta$ equal $z$ on sources and $t$ on sinks, and write
\[
 p=\theta_P,\ q=\theta_Q,\ r=\theta_R,\ d_0=\theta_D,
 \ e=\theta_E,\ f=\theta_F,\ g=\theta_G,\ l=\theta_L,
 \ m=\theta_M,\ n=\theta_N,\quad
 \alpha=\theta_A,\ \beta=\theta_B.
\]
These are proof variables only; in particular $p\ge1$.

Whenever $a(S)=b(R)$, both $(S,R)$ and its complement have worth
equal to their sink capacity, and their worths sum to $196$.
Efficiency makes both core inequalities tight, so
$z(S)-t(R)=\mathbf1_{B\in S}$. Equivalently, if a subset $W$ of all
players has baseline weight $98$, then
$\theta(W)=\sum_jt_j+\mathbf1_{B\in W}$.
For every $C_0$ subset $K$ of baseline weight $55$, apply this to
$W=\{A,B\}\cup K$ and use~\eqref{flow:loss-total} to obtain
\[
                     \theta(K)=k_0:=p+g+m.
\]
The eight such subsets
\[
 PGM,\ PGN,\ PEFM,\ PQDFM,\ RDEM,\ PQGL,\ PQREFG,\ PDEFG
\]
therefore have equal $\theta$-sum. In succession, these equalities give
\[
 m=n,\quad g=e+f,\quad e=q+d_0,\quad p+f=r+d_0,
 \quad m=q+l,\quad m=q+r+g,\quad m=d_0+g.
\]
Solving them yields
\begin{align}
 d_0&=q+r,&e&=2q+r,&f&=q+2r-p,&g&=3q+3r-p,\notag\\
 l&=3q+4r-p,&m=n&=4q+4r-p,&k_0&=7q+7r-p.
       \label{flow:identities}
\end{align}
Thus $\theta(C_0)=-4p+19q+20r$. Combining this with efficiency gives
\begin{equation}\label{flow:loss-formulas}
 \alpha+\beta=1-2p+5q+6r,\qquad
                 \sum_i z_i=1-3p+12q+13r.
\end{equation}

If $W$ has baseline weight $97$, take $S=W\cap\mathrm{sources}$
and $R=\mathrm{sinks}\setminus W$. Then $a(S)-b(R)=-1$, so its
actual source capacity is strictly smaller than its sink capacity.
Core coverage gives $z(S)\le t(R)$, or
$\theta(W)\le\sum_jt_j$. Hence for every $C_0$ subset $K$ of
weight $54$,
\[
                         \theta(K)\le k_0-1.
\]
Use $K_1=\{M,E,D,Q,P\}$ and $K_2=\{L,F,D,R,P\}$.
Their $\theta$-sums from~\eqref{flow:identities} are $8q+6r$ and
$5q+8r-p$, respectively, so
\[
              r-q\ge p+1,\qquad r\le2q-1.
\]
Consequently
\begin{equation}\label{flow:qr}
                     q\ge p+2,\qquad r\ge2p+3,\qquad p\ge1.
\end{equation}
For the weight-$97$ subset $\{B,M,L,P\}$, whose $C_0$ portion has
$\theta$-sum $k_0+r$, the same strict-coalition inequality gives
$\alpha\ge r+1\ge6$. For the weight-$97$ subset
$\{A\}\cup(C_0\setminus\{M,N\})$, it gives
\[
                    \beta\ge -p+4q+5r+1\ge37.
\]
Also~\eqref{flow:identities} and~\eqref{flow:qr} imply
$g\ge8p+15\ge23$ and $m\ge11p+20\ge31$.
Thus every source payoff is coordinatewise bounded above by
$(1,6,23,31,38)$. Vector $U$ attains the first four bounds and $V$
attains the fifth. For each sink, deletion gives $H_j\le c_j$;
the all-sink cut allocation attains all those bounds at once.

Finally~\eqref{flow:loss-formulas}--\eqref{flow:qr} give
\[
             \sum_i z_i\ge35p+64\ge99,
              \qquad \alpha+\beta\ge15p+29\ge44.
\]
Source total is therefore at most $197-99=98$, and $x_A+x_B$ at
most $12+75-44=43$. Both $U,V$ attain those support bounds. This
proves the claimed incompatibility. Vector $U$ already dominates
$H/3$ on every source and sink and is efficient, so the protected
subsidy is zero.
\end{proof}

The baseline pattern is drawn from Isbell's integer majority-weight
representations~\cite[p.~27]{flow:Isbell1959}; the weights are not
claimed as a new construction ingredient. Integer nonuniqueness alone
would not prove Theorem~\ref{flow:endpoint}, which uses additional
all-coalition and real-valued core arguments. With all allowed-pair
returns one, this is also unrestricted unweighted transportation.
It matches the historical side-optimal-core question stated in
N\'u\~nez's 2013 lecture slides~\cite[pp.~37--40]{flow:Nunez2013}.
Neither its novelty relative to later or unread literature nor the
current openness of that historical question is established here.
The weighted multiple-partners model with capacity one per pair is
different and is not settled by this example.

\subsection{Why mixed tight-set uncrossing is insufficient}
For endpoint sources $I$, sinks $J$, and fixed public middle capacities,
let $w(Q,R)$ be the flow value when endpoint sets $Q\subseteq I$ and
$R\subseteq J$ are active. The established network-rank inequality is
\begin{equation}\label{flow:mixed}
 w(Q,R)+w(Q',R')\ge
 w(Q\cap Q',R\cup R')+w(Q\cup Q',R\cap R').
\end{equation}
This is the mixed submodularity of network linking
ranks~\cite{flow:Goemans2011}. A direct cut proof also fixes the direction
needed here. Use the full fixed vertex set, minimum-cut source sides
$U,V$ for the original two endpoint profiles, and source sides
$U\cap V,U\cup V$ for the two crossed profiles in the order displayed.
The public-middle cut contributions satisfy submodularity. For a source
arc, write $a,a'$ for its activity indicators and $u,v$ for the
indicators that its head is outside $U,V$. Then
\[
 au+a'v\ge (a\wedge a')\max(u,v)+(a\vee a')\min(u,v).
\]
For a sink arc use activity indicators $b,b'$ and tail-membership
indicators $u,v$; its crossed contribution is
$(b\vee b')\min(u,v)+(b\wedge b')\max(u,v)$, also at most
$bu+b'v$. Multiply by nonnegative capacities, add the middle-arc
inequality, and bound each crossed cut below by its maximum-flow value.
This proves~\eqref{flow:mixed}.

Additive payoffs have equal total payment on an original pair and its
mixed-crossed pair. The inequality~\eqref{flow:mixed} has the wrong
direction to force the crossed coalitions tight when the originals are
tight. The failure occurs even in a two-hop game: two source owners have
capacities $3,3$, three sink owners capacities $2,2,1$, and the core
allocation is $(0,0;2,2,1)$. The coalitions
\[
 (\{a_1\},\{b_1,b_3\}),\qquad
 (\{a_1,a_2\},\{b_2,b_3\})
\]
each have worth and payment $3$. Their crossed pair has worths $3,1$
but payments $5,1$. Thus total worth falls from $6$ to $4$ while
payment stays $6$, although the original coalitions overlap and neither
contains the other. This disproves mixed-crossed tight-set closure, not
every conceivable use of mixed submodularity or side-optimality.

\paragraph{Remaining scope.}
The general rational-capacity conjecture
$s_{1/3}(G)\le T/3$ for all arc-player networks of path length at most
three remains unresolved by these results. The endpoint example has
subsidy zero, and the unit and quota theorems make restricted assumptions.
For example, the serial path of capacities $(2,2,1)$ has the unit-simplex
entire core and $H=(1,1,1)$, although only its last arc lies in a minimum
cut. Clipping that cut allocation at $H/3$ costs $2/3$, whereas the
protected subsidy is zero at $(1/3,1/3,1/3)$. This illustrates why
entire-core maxima cannot be replaced by cut or optimal-dual maxima at
unequal capacities. No owner cloning, simple-graph transfer, or all-profile
truthfulness assertion is implicit in this appendix.
